\subsection{分裂复形}

命题6. — 设 (C, d) 是一个复形。以下条件等价：

\begin{enumerate}
\def\labelenumi{(\roman{enumi})}
\item
  存在从 (C, d) 到 (H(C), 0) 的同伦等价
\item
  存在一个次数为1的A-分次自同态s，使得 \(d = d \circ s \circ d\) ;
\item
  B(C) 与 Z(C) 是 C 的直和项；
\item
  (C, d) 是子复形的直和，这些子复形要么长度为 0，要么长度为 1 且同调为零。
\item
  \(\Rightarrow\) (ii)：设 \(\varphi: \mathbf{C} \to \mathbf{H}(\mathbf{C})\) 是一个同伦等价；则存在一个复形态射 \(\psi: \mathbf{H}(\mathbf{C}) \to \mathbf{C}\) 以及一个自同态 \(s\) （ \(\mathbf{C}\) 的、次数为 1 的），使得 \(\psi \circ \varphi = 1_{\mathbf{C}} - s \circ d - d \circ s\) 。我们有 \(d \circ \psi = \psi \circ 0 = 0\) ，因此
\end{enumerate}

\[
0 = d \circ \psi \circ \varphi = d - d \circ s \circ d - d \circ d \circ s = d - d \circ s \circ d, \quad \text {whence (ii)}.
\]

\begin{enumerate}
\def\labelenumi{(\roman{enumi})}
\setcounter{enumi}{1}
\item
  \(\Rightarrow\) (iii)：设 \(s\) 如(ii)中所述。则 \(d \circ (1_{\mathbb{C}} - s \circ d) = 0\) ，因此 \(1_{\mathbb{C}} - s \circ d\) 是 \(\mathbf{C}\) 到 \(\mathbf{Z}(\mathbf{C})\) 上的投影映射，且 \((d \circ s) \circ d = d\) ，因此 \(d \circ s\) 是 \(\mathbf{C}\) 到 \(\mathbf{B}(\mathbf{C})\) 上的投影映射 .
\item
  \(\Rightarrow\) (iv)：对每个 \(n\in\mathbf{Z}\) ，设 \(Z_{n}=Z_{n}(\mathrm{C})\) , \(B_{n}=B_{n}(\mathrm{C})\) ，并选取子模 \(K_{n}\) 与 \(B_{n}^{\prime}\) （含于 \(C_{n}\) ），使得 \(C_{n}=B_{n}^{\prime}\oplus Z_{n}\) , \(Z_{n}=K_{n}\oplus B_{n}\) 。然后
\end{enumerate}

\[
\mathrm{E} _ {(n)} = \mathrm{K} _ {n} (- n) \quad \text { and } \quad \mathrm{F} _ {(n)} = \mathrm{B} _ {n} ^ {\prime} (- n) \oplus \mathrm{B} _ {n - 1} (1 - n)
\]

是 (C, d) 的子复形；我们有 (C, d) = \(\bigoplus_{n \in \mathbf{Z}} (\mathrm{E}_{(n)} \oplus \mathrm{F}_{(n)})\) ；每个 \(\mathrm{E}_{(n)}\) 要么为零，要么长度为 0，每个 \(\mathrm{F}_{(n)}\) 要么为零，要么长度为 1 且同调为零，由此得 (iv)。

\begin{enumerate}
\def\labelenumi{(\roman{enumi})}
\setcounter{enumi}{3}
\tightlist
\item
  \(\Rightarrow\) (i)：只需注意当 C 的长度为零，或同调为零且长度为 1 时，(i) 成立。
\end{enumerate}

定义 6. --- 若复形 C 满足命题 6 的等价条件，则称其为分裂的。

满足命题 6 条件 (ii) 的 C 的自同态 s 称为 C 的一个分裂。

例子. --- 1) 具有零微分的复形是分裂的。

\begin{enumerate}
\def\labelenumi{\arabic{enumi})}
\setcounter{enumi}{1}
\item
  同伦于零的复形正是同调为零的分裂复形，即满足 H(C) = 0 且 Z(C) 是 C 的直和项的复形 C。
\item
  设 \(f: \mathbf{M} \to \mathbf{N}\) 是 A-模的同态，且 C 是满足 \(\mathbf{C}_1 = \mathbf{M}, \mathbf{C}_0 = \mathbf{N}, \mathbf{C}_i = 0\) 当 \(i \neq 0, 1, d_1 = f, d_i = 0\) 当 \(i \neq 1\) 的复形。则 C 是分裂的当且仅当 Ker \(f\) 是 M 的直和项且 Im \(f\) 是 N 的直和项。
\item
  只要 B(C) 与 H(C) 是投射的（或只要对每个 n， \(B_{n}(C)\) 与 \(H_{n}(C)\) 是内射的），复形 C 就是分裂的。事实上，由第 1 号的正合列 \((I_{n})\) 到 \((IV_{n})\) 可知，此时 Z(C) 是 C 的直和项且 B(C) 是 Z(C) 的直和项（相应地，B(C) 是 C 的直和项且 Z(C)/B(C) 是 C/B(C) 的直和项）。
\item
  特别地，若 A 是主理想环，则自由复形 C 是分裂的当且仅当 H(C) 是自由的（即对每个 n ∈ Z，H \(_{n}\) (C) 是自由的）。
\end{enumerate}

注. --- a) 假设分次 A-模的典范正合列

\[
0 \to \mathrm{B} (\mathrm{C}) \to \mathrm{Z} (\mathrm{C}) \xrightarrow {\pi} \mathrm{H} (\mathrm{C}) \to 0\tag{II}
\]

是分裂的（例如当 H(C) 是投射的，或对每个 n，B \(_{n}\) (C) 是内射的时成立）；设 σ : H(C) → Z(C) 是 π 的一个分次 A-线性截面，并设 ψ 为从 H(C) 到 C 的映射 x ↦ σ(x)。则 ψ 是从 (H(C), 0) 到 C 的一个拟同构，使得 H(ψ) = 1 \(_{H(C)}\) .

\begin{enumerate}
\def\labelenumi{\alph{enumi})}
\setcounter{enumi}{1}
\tightlist
\item
  假设分次 A-模的典范正合列
\end{enumerate}

\[
0 \rightarrow \mathrm{H} (\mathrm{C}) \stackrel {i} {\rightarrow} \mathrm{C} / \mathrm{B} (\mathrm{C}) \rightarrow \mathrm{B} (\mathrm{C}) (- 1) \rightarrow 0\tag{IV}
\]

是分裂的（例如当 B(C) 是投射的，或 H \(_{n}\) (C) 对每个 \(n\) 是内射的时成立）；设 \(\tau: \mathbf{C} / \mathbf{B}(\mathbf{C}) \to \mathbf{H}(\mathbf{C})\) 是 \(i\) 的一个分次 A-线性收缩，并设 \(\varphi\) 为从 C 到 \(\mathbf{H}(\mathbf{C})\) 的同态，它把 C 的每个元素映到该元素在 \(\tau\) 下的像，即该元素模 \(\mathbf{B}(\mathbf{C})\) 的类的像。则 \(\varphi\) 是从 C 到 \((\mathbf{H}(\mathbf{C}), 0)\) 的一个拟同构，使得 \(\mathbf{H}(\varphi) = 1_{\mathbf{H}(\mathbf{C})}\) .

